
\begin{enumerate}[label=\thesubsection.\arabic*, ref=\thesubsection.\theenumi]
\item
Let $ax^2 + bx + c =0$ be a quadratic equation.Let $\alpha ,\beta$ be the roots of the equation.
\begin{align}
\text{The roots of the equation are} \frac{-b \pm \sqrt{b^2 - 4ac}}{2a}\label{eq:EE26BTECH11211-Prob_1.1.10}
\end{align}
\begin{align}
\text{sum of the roots of the equation }= \alpha + \beta = -\frac{b}{a}\label{eq:EE26BTECH11211-Prob_1.2.10}
\end{align}
\begin{align}
\text{product of the roots of the equation}= \alpha\beta=\frac{c}{a}\label{eq:EE26BTECH11211-Prob_1.3.10}
\end{align}
\begin{align}
\text{Discriminant }= D =b^2 - 4ac.\label{eq:EE26BTECH11211-Prob_1.4.10}
\end{align}
\\$b^2 - 4ac = 0$, if roots are real and equal. 
\\$b^2 - 4ac > 0$, if roots are real and distinct.
\\$b^2 - 4ac < 0$, if roots are complex.

\item
Let $a,a_2,a_3,.....,a_n$ be the n terms of arithmetic progression.Then,
\begin{align}
\text{common difference d}=a_n - a_{n-1}
\end{align}
\begin{align}
\text{General equation for nth term of A.P series is } a_n = a + \brak{n-1}d\label{eq:EE26BTECH11211-Prob_2.2.10}
\end{align}
\begin{align}
\text{Sum of first n terms } S_n = \frac{n}{2}\sbrak{2a + \brak{n-1}d}\label{eq:EE26BTECH11211-Prob_2.3.10}
\end{align}
\begin{align}
\text{Sum of first n terms } S_n = \frac{n}{2}\sbrak{a + l} \label{eq:EE26BTECH11211-Prob_2.4.10}
\end{align}
\item
Let $A\brak{x_1,y_1}$ and $B\brak{x_2,y_2}$ be two points then the distance between two points A and B is.
\begin{align}
AB = \sqrt{\brak{x_2 - x_1}^2 + \brak{y_2 - y_1}^2} \label{eq:EE26BTECH11211-Prob_3.10}
\end{align}

\item
Let $x_1,x_2,....,x_n$ be some data and $f_1,f_2,...,f_n$ be their respective frequencies then the mean of the data is 
\begin{align}
\bar{x} =\frac{\sum_{i=1}^{n} f_i x_i}{\sum_{i=1}^{n} f_i}
\end{align}

\item
Surface area of sphere
\begin{align}
S = 4\pi r^2
\end{align}

\item
Curved surface area of hemisphere
\begin{align}
S = 2\pi r^2
\end{align}

\item
Total surface area of a hemisphere
\begin{align}
S = 3\pi r^2
\end{align}

\item
Volume of the cone.
\begin{align}
V = \frac{1}{3}\pi r^2 h
\end{align}

\item 
Let $A =\myvec{a & b\\ c & d}$ then the determinant of matrix A is.
\begin{align}
\mydet{A} = \mydet{a & b\\ c & d} = ad - bc
\end{align}

\item
Let $A =\myvec{a & b\\ c & d}$ then the inverse of matrix A exists only when $\mydet{A}\neq{0}$ and the inverse is,
\begin{align}
A^{-1} = \frac{1}{ad - bc}\myvec{d & -b\\ -c & a}\label{eq:EE26BTECH11211_Prob_3.10.10}
\end{align}

\item
$sin\brak{A + B} = sinAcosB + cosAsinB$

\item
Let $P\brak{x ,y}$ divides the line segment joining $A\brak{x_1,y_1}$ and $B\brak{x_2,y_2}$ in the ratio m:n.
\begin{align}
\brak{x,y} = \brak{\frac{mx_2 + nx_1}{m + n},\frac{my_2 + ny_1}{m + n}} \label{eq:EE26BTECH11211_Prob_3.12.10}
\end{align}

\item
Volume of the cylinder is
\begin{align}
V = \pi r^2 h
\end{align}

\item
Lenth of the arc of the circle subtending an angle $\theta$ is
\begin{align}
l = \frac{\theta}{360^\circ}2\pi r \label{eq:EE26BTECH11211_Prob_3.14.10}
\end{align} 

\item
Area of the arc of the circle subtending an angle $\theta$ is
\begin{align}
l = \frac{\theta}{360^\circ}\pi r^2 \label{eq:EE26BTECH11211_Prob_3.15.10}
\end{align} 

\item
The formula for the median of grouped data is
\begin{align}
Median = l +\brak{\frac{\frac{N}{2} - cf}{f}}h \label{eq:EE26BTECH11211_Prob_3.16.10}
\end{align} 
N= total frequency
\\cf = cummulative frequency of class preceding the median class
\\h = class size
\\l = lower limit of the median class

\item
The formula for mode of grouped data is 
\begin{align}
Mode = l + \brak{\frac{f_1 - f_0}{2f_1 - f_0 -f_2}}h \label{eq:EE26BTECH11211_Prob_3.17.10}
\end{align}
$f_1 $= frequency of the modal class
\\$f_0$ = frequency of class preceding the modal class
\\$f_2$ = frequency of the class succeeding the modal class
\\h = class size
\\l = lower limit of the modal class

\end{enumerate}
